\section{GANs com Operações Convolucionais}

\textbf{Questão 2:} Escolha uma base de dados de imagens e treine um modelo GAN sobre ela, com um gerador utilizando operações convolucionais.

\begin{itemize}
    \item[(a)] Considere dois critérios de parada diferentes e compare-os;
    \item[(b)] Escolha duas métricas de qualidade de imagem e compare-as para analisar o resultado da GAN.
\end{itemize}

\textbf{Resolução:}

Para a realização deste experimento, foi utilizado o conjunto de dados \textit{MNIST} (\textit{Modified National Institute of Standards and Technology}) \cite{lecun1998gradient}, composto por imagens em escala de cinza de dígitos manuscritos. Cada imagem possui dimensão de $28 \times 28$ pixels e pertence a uma das dez classes correspondentes aos dígitos de $0$ a $9$. O conjunto de dados contém 70.000 imagens, sendo 60.000 destinadas ao treinamento e 10.000 ao teste. Neste trabalho, apenas o conjunto de treinamento foi utilizado para o treinamento da Rede Generativa Adversarial (GAN). A Figura~\ref{fig:data2} apresenta alguns exemplos de imagens presentes na base de dados. Além disso, através da Figura~\ref{fig:data_distribution2}, podemos observar que o conjunto de dados é aproximadamente balanceado em relação ao número de amostras de cada classe disponíveis para treinamento.

Como a camada de saída da rede geradora utilizará a função de ativação \textit{Tanh}, cujos valores pertencem ao intervalo $[-1,1]$, as imagens do MNIST, originalmente representadas no intervalo $[0,1]$, foram normalizadas para o intervalo $[-1,1]$ por meio da seguinte transformação:

\begin{equation}
x_{\text{norm}} = 2x - 1.
\end{equation}

\begin{figure}[H]
    \centering
    \includegraphics[width=0.5\textwidth]{imgs/data2.png}
    \caption{Exemplos de imagens do conjunto de treinamento do MNIST.}
    \label{fig:data2}
\end{figure}

\begin{figure}[H]
    \centering
    \includegraphics[width=0.75\textwidth]{imgs/mnist_distribution2.png}
    \caption{Distribuição dos dados de treinamento por classes.}
    \label{fig:data_distribution2}
\end{figure}

Considerando que o problema envolve a geração de imagens, foi adotada uma arquitetura convolucional para a GAN, inspirada na estrutura proposta pelas DCGANs (\textit{Deep Convolutional Generative Adversarial Networks}). A arquitetura utilizada para as redes é ilustrada na Figura~\ref{fig:gan2}.

Como podemos observarm, a rede geradora recebe como entrada um vetor de ruído aleatório de dimensão 100, representado no formato $(100,1,1)$. A partir desse vetor, sucessivas camadas de convolução transposta (\textit{ConvTranspose2d}) realizam o aumento gradual da resolução espacial até produzir uma imagem sintética de dimensão $28 \times 28$. Entre as camadas convolucionais são empregadas camadas de \textit{Batch Normalization} e a função de ativação ReLU, enquanto a camada de saída utiliza a função \textit{Tanh}, garantindo que os valores gerados permaneçam no mesmo intervalo dos dados normalizados.

Por outro lado, a rede discriminadora recebe como entrada imagens reais ou sintéticas de dimensão $28 \times 28$ e realiza sucessivas operações de convolução para reduzir sua dimensão espacial e extrair características discriminativas. Nas camadas intermediárias é utilizada a função de ativação \textit{LeakyReLU} com coeficiente negativo igual a 0,2, além de uma camada de \textit{Batch Normalization}. A camada final produz um único valor escalar, indicando o grau de confiança da rede de que a imagem fornecida pertence ao conjunto de dados reais.

A Tabela~\ref{tab:gan_architecture} resume a arquitetura completa das redes geradora e discriminadora implementadas neste trabalho. Os hiperparâmetros da arquitetura foram definidos empiricamente, a partir de experimentos envolvendo diferentes configurações de número de filtros, tamanhos de kernel e demais parâmetros das camadas convolucionais.

\begin{table}[H]
\centering
\caption{Arquitetura da GAN utilizada nos experimentos.}
\label{tab:gan_architecture}
\renewcommand{\arraystretch}{1.2}
\begin{tabular}{|c|c|c|c|c|c|}
\hline
\textbf{Rede} & \textbf{Camada} & \textbf{Filtros} & \textbf{Kernel} & \textbf{Stride} & \textbf{Padding} \\
\hline
\multirow{3}{*}{Geradora}
& ConvTranspose2d & 256 & $7\times7$ & 1 & 0 \\
& ConvTranspose2d & 128 & $4\times4$ & 2 & 1 \\
& ConvTranspose2d & 1 & $4\times4$ & 2 & 1 \\
\hline
\multirow{3}{*}{Discriminadora}
& Conv2d & 64 & $4\times4$ & 2 & 1 \\
& Conv2d & 128 & $4\times4$ & 2 & 1 \\
& Conv2d & 1 & $7\times7$ & 1 & 0 \\
\hline
\end{tabular}
\end{table}

Na rede geradora, as duas primeiras camadas convolucionais transpostas são seguidas por operações de \textit{Batch Normalization} e pela função de ativação ReLU. Já na camada de saída é utilizada a função \textit{Tanh}. Na rede discriminadora, a primeira camada convolucional é seguida por uma função \textit{LeakyReLU}; a segunda camada utiliza \textit{Batch Normalization} e \textit{LeakyReLU}; enquanto a camada final não emprega função de ativação, produzindo diretamente o escore utilizado pela função de perda durante o treinamento.

\begin{figure}[H]
    \centering
    \includegraphics[width=0.95\textwidth]{imgs/gan2.png}
    \caption{Arquitetura da GAN convolucional utilizada neste trabalho.}
    \label{fig:gan2}
\end{figure}

O treinamento foi realizado de maneira análoga ao descrito na Questão~\ref{cap:exercise1}, utilizando o otimizador Adam com os mesmos hiperparâmetros e as mesmas funções de custo para o gerador e o discriminador.

Considerando um limite máximo de 100 épocas, foram avaliadas duas estratégias distintas de parada antecipada, descritas a seguir.

\begin{itemize}
    \item \textbf{(M1) Early Stopping baseado no FID:} nesta estratégia, a qualidade das imagens geradas foi monitorada por meio da métrica \textit{Fréchet Inception Distance} (FID), definida por
    \begin{equation}
        \mathrm{FID}(\mathbf{x},\mathbf{g}) =
        \left\|\boldsymbol{\mu}_x-\boldsymbol{\mu}_g\right\|_2^2 +
        \operatorname{Tr}
        \left(
        \Sigma_x+\Sigma_g-
        2(\Sigma_x\Sigma_g)^{1/2}
        \right),
        \label{eq:fid}
    \end{equation}
    onde $(\boldsymbol{\mu}_x,\Sigma_x)$ representam a média e a matriz de covariância das características extraídas das imagens reais e $(\boldsymbol{\mu}_g,\Sigma_g)$ representam as estatísticas correspondentes das imagens geradas. Quanto menor o valor do FID, maior a similaridade entre as distribuições das imagens reais e sintéticas.

    O treinamento era interrompido caso o melhor valor de FID não apresentasse melhoria durante uma janela consecutiva de $n=10$ épocas.

    \item \textbf{(M2) Estabilidade da acurácia do discriminador:} nesta estratégia, foi monitorada a acurácia do discriminador sobre o conjunto de validação ao final de cada época. A hipótese adotada é que, quando essa acurácia passa a oscilar dentro de uma faixa reduzida por um período suficientemente longo, o treinamento pode ter atingido um estado próximo do equilíbrio entre gerador e discriminador.

    Para verificar essa condição, foi considerada uma janela móvel contendo as últimas 10 épocas. O treinamento era encerrado quando a diferença entre a maior e a menor acurácia observadas nessa janela fosse inferior a $10\%$, caracterizando um comportamento suficientemente estável do discriminador.
\end{itemize}

Após definir os critérios de parada, a GAN foi treinada utilizando cada uma das estratégias propostas. A Figura~\ref{fig:loss_gan2} apresenta a evolução das funções de custo normalizadas do gerador e do discriminador durante o treinamento. Observa-se que a estratégia M1 interrompeu o treinamento na época 26, momento em que o critério de estagnação do FID foi satisfeito. Em contrapartida, utilizando a estratégia M2, o treinamento prosseguiu até a época 59, quando foi detectada a estabilidade da acurácia do discriminador.

\begin{figure}[H]
    \centering
    \includegraphics[width=0.60\textwidth]{imgs/loss_gan2.png}
    \caption{Evolução das funções de custo normalizadas do gerador e do discriminador durante o treinamento da GAN para as duas estratégias de parada antecipada avaliadas: Early FID (M1) e estabilidade da acurácia do discriminador (M2). Observa-se que o treinamento utilizando M1 foi interrompido na época 26, enquanto o treinamento utilizando M2 prosseguiu até a época 59.}
    \label{fig:loss_gan2}
\end{figure}

Para analisar o comportamento de ambos os critérios, a Figura~\ref{fig:FID2} apresenta a evolução do FID ao longo das épocas, enquanto a Figura~\ref{fig:disc_acc2} mostra a evolução da acurácia do discriminador.

\begin{figure}[H]
    \centering
    \includegraphics[width=0.60\textwidth]{imgs/FID2.png}
    \caption{Evolução da métrica Fréchet Inception Distance (FID) ao longo do treinamento para as estratégias Early FID (M1) e estabilidade da acurácia do discriminador (M2). Valores menores indicam maior similaridade entre as distribuições das imagens reais e geradas.}
    \label{fig:FID2}
\end{figure}

\begin{figure}[H]
    \centering
    \includegraphics[width=0.60\textwidth]{imgs/disc_acc2.png}
    \caption{Evolução da acurácia do discriminador no conjunto de validação durante o treinamento da GAN para as estratégias Early FID (M1) e estabilidade da acurácia do discriminador (M2). A região final da curva de M2 apresenta menor variação entre épocas, satisfazendo o critério de estabilidade utilizado para a parada antecipada.}
    \label{fig:disc_acc2}
\end{figure}

Observando a Figura~\ref{fig:FID2}, nota-se que a estratégia M1 interrompeu o treinamento quando o FID permaneceu sem melhora durante a janela de monitoramento estabelecida. Entretanto, ao analisar a curva correspondente ao treinamento utilizando M2, percebe-se que, após algumas épocas adicionais, o FID voltou a diminuir e atingiu valores inferiores ao obtido pela estratégia M1. Esse comportamento indica que, para os hiperparâmetros utilizados, o critério de parada baseado exclusivamente na estagnação do FID mostrou-se excessivamente conservador, interrompendo o treinamento antes que o gerador atingisse sua melhor qualidade.

A Figura~\ref{fig:disc_acc2} permite interpretar o comportamento da estratégia M2. Entre aproximadamente as épocas 30 e 50, a acurácia do discriminador apresentou oscilações relativamente elevadas, indicando que o processo adversarial ainda estava em evolução. Somente após esse período a acurácia passou a permanecer dentro da faixa de estabilidade estabelecida, levando ao encerramento do treinamento. Embora os valores absolutos da acurácia nesse momento sejam relativamente elevados, sugerindo que o discriminador ainda conseguia distinguir parte das imagens sintéticas, a evolução do FID demonstra que o gerador continuou melhorando sua capacidade de produzir amostras realistas durante esse intervalo.

A Figura~\ref{fig:res2} apresenta exemplos de imagens sintetizadas pelos modelos obtidos através das duas estratégias de parada. Visualmente, observa-se uma ligeira superioridade das imagens produzidas pela estratégia M2, que apresentam menor quantidade de artefatos no fundo e dígitos com formas mais consistentes.

\begin{figure}[H]
    \centering

    \begin{subfigure}[b]{0.48\textwidth}
        \centering
        \includegraphics[width=0.85\textwidth]{imgs/res21.png}
        \caption{Imagens geradas pelo modelo treinado utilizando o critério de parada antecipada baseado no FID (M1).\\}
        \label{fig:res21}
    \end{subfigure}
    \hfill
    \begin{subfigure}[b]{0.48\textwidth}
        \centering
        \includegraphics[width=0.85\textwidth]{imgs/res22.png}
        \caption{Imagens geradas pelo modelo treinado utilizando o critério de estabilidade da acurácia do discriminador (M2).}
        \label{fig:res22}
    \end{subfigure}

    \caption{Exemplos de dígitos sintéticos produzidos pelos modelos obtidos utilizando os dois critérios de parada antecipada. Visualmente, observa-se que o modelo treinado com o critério M2 produz amostras mais consistentes, com menor quantidade de artefatos e melhor definição dos dígitos.}
    \label{fig:res2}
\end{figure}

Como referência adicional, o modelo também foi treinado durante as 100 épocas previstas, sem a utilização de parada antecipada. A Tabela~\ref{tab:fid_results} resume os resultados quantitativos obtidos utilizando as métricas FID e MMD para os três modelos treinados.

\begin{table}[H]
\centering
\caption{Comparação entre os diferentes critérios de parada.}
\label{tab:fid_results}
\begin{tabular}{lcc}
\hline
\textbf{Estratégia} & \textbf{Épocas} & \textbf{FID} $\downarrow$ \\
\hline
Treinamento completo & 100 & 0.1075\\
Early FID (M1) & 26 & 0.1164\\
Estabilidade do discriminador (M2) & 59 & \textbf{0.0727}\\
\hline
\end{tabular}
\label{tab:fid_results}
\end{table}

Os resultados quantitativos corroboram a análise visual. Apesar de o treinamento completo utilizar todas as 100 épocas disponíveis, ele não apresentou o melhor desempenho. A estratégia M2 obteve simultaneamente o menor valor de FID e de MMD, indicando maior proximidade entre as distribuições das imagens reais e geradas. Por outro lado, a estratégia M1 apresentou o pior desempenho entre os três cenários, reforçando que, para esta configuração experimental, o critério de parada baseado apenas na estagnação do FID interrompeu o treinamento prematuramente.

De maneira geral, os experimentos indicam que a utilização da estabilidade da acurácia do discriminador como critério de parada mostrou-se uma alternativa mais eficiente para este problema, permitindo que o gerador continuasse evoluindo mesmo após períodos temporários de estagnação do FID e produzindo, ao final do treinamento, amostras de melhor qualidade.